% Chapter 8: Optimal Statistical Tests
% Hogg, McKean, Craig - Introduction to Mathematical Statistics

\chapter{Optimal Statistical Tests}\label{chap:optimal-tests}

\section*{Chapter 8 Overview}

在前几章中，我们建立了估计理论和假设检验的基本框架。我们学会了如何构造点估计、区间估计，以及如何使用似然比检验。但在所有这些讨论中，一个根本性的问题始终潜伏在幕后：

\begin{center}
\textit{在所有可能的检验中，哪个检验是``最好''的？}
\end{center}

这个问题——即\textbf{最优检验理论（optimal test theory）}——是本章的核心。在假设检验中，我们面临一个固有的权衡（trade-off）：降低第一类错误（Type I error）的概率往往意味着提高第二类错误（Type II error）的概率，反之亦然。最优检验理论试图在固定的显著性水平下，找到功效（power）最高的检验。

\textbf{为什么需要最优性理论？} 考虑一个医学检测场景：我们希望设计一个检验来筛查某种疾病。给定显著性水平 $\alpha = 0.05$，我们当然希望当患者真的患病时，检验能以最高概率检测出来——这就是功效。没有最优性理论，我们只能在黑暗中摸索不同的检验方法，无法保证我们选择的检验是最佳的。

本章建立这套理论的路径如下：

\begin{enumerate}
\item \textbf{Neyman-Pearson 引理}（8.1 节）：对于简单假设对简单假设的情况，给出了最优检验的完整刻画——这是整个理论的基石。
\item \textbf{一致最大功效检验（UMP）}（8.2 节）：将最优性概念推广到复合备择假设。介绍 monotone likelihood ratio（MLR）和 Karlin-Rubin 定理——单参数指数族存在 UMP 检验。
\item \textbf{似然比检验（LRT）}（8.3 节）：在一般复合假设下，LRT 提供了一种构造检验的通用方法。虽然 LRT 不一定是最优的，但在很多情况下它是渐近最优的。
\item \textbf{序贯概率比检验（SPRT）}（8.4 节）：打破固定样本量的束缚，允许检验在数据积累过程中随时停止——Wald 的革命性贡献。
\item \textbf{极小化极大分类}（8.5 节）：将决策论框架应用于假设检验和分类问题。
\end{enumerate}

\section{8.1 Most Powerful Tests}\label{sec:8.1}

\subsection{问题的起源：什么是"最好"的检验}

在 4.5 节中，我们引入了假设检验的基本框架：给定原假设 $H_0$ 和备择假设 $H_1$，我们构造一个临界区域 $C$，当样本落入 $C$ 时拒绝 $H_0$。但正如我们在 4.5 节所见，对于同一对假设，存在无穷多个可能的检验——临界区域 $C$ 的选择并非唯一。

\begin{center}
\textit{在所有这些检验中，我们凭什么说一个比另一个更好？}
\end{center}

要回答这个问题，我们首先需要一个评价标准。在假设检验中，这个标准就是\textbf{功效（power）}——当 $H_1$ 为真时拒绝 $H_0$ 的概率。直观上，我们希望：

\begin{enumerate}
\item 在 $H_0$ 为真时，拒绝 $H_0$ 的概率（即 $\alpha$，显著性水平）不超过预设值；
\item 在 $H_1$ 为真时，拒绝 $H_0$ 的概率（即功效）尽可能高。
\end{enumerate}

这就引出了"最好检验"的定义。

\subsection{简单假设对简单假设的最优检验}

设 $X_1, \ldots, X_n$ 是来自分布 $f(x;\theta)$ 的随机样本，其中 $\theta \in \Omega = \{\theta', \theta''\}$（即只有两个可能的参数值）。我们考虑检验：

\[
H_0: \theta = \theta' \quad \text{versus} \quad H_1: \theta = \theta''. \tag{8.1.1}\label{eq:simple-vs-simple}

\]

\begin{definition}[最佳临界区域]\label{def:best-critical-region}
设 $C$ 是样本空间的一个子集。如果
\begin{enumerate}
\item[(a)] $P_{\theta'}[\mathbf{X} \in C] = \alpha$（大小为 $\alpha$）；
\item[(b)] 对每一个满足 $P_{\theta'}[\mathbf{X} \in A] = \alpha$ 的子集 $A$，都有
\[
P_{\theta''}[\mathbf{X} \in C] \geq P_{\theta''}[\mathbf{X} \in A],
\]
\end{enumerate}
则称 $C$ 是检验 $H_0: \theta = \theta'$ 对 $H_1: \theta = \theta''$ 的\textbf{大小为 $\alpha$ 的最佳临界区域（best critical region of size $\alpha$）}。
\end{definition}

\textbf{定义的直观含义}：在所有满足显著性水平 $\alpha$ 的检验中，最佳临界区域对应的检验在 $H_1$ 为真时具有最高的拒绝概率——这正是我们想要的。

\textbf{为什么要关心最佳检验？} 设想你在设计一个药物试验。你设定 $\alpha = 0.05$（即当药物无效时，只有 5\% 的概率错误地拒绝 $H_0$）。在所有满足这个条件的检验中，你当然希望当药物真的有效时，以最高概率检测出来——这就是最佳检验的价值。

那么，这样的最佳检验是否存在？Neyman-Pearson 定理给出了答案。

\subsection{Neyman-Pearson 定理}

\begin{center}
\textit{有没有一种系统性的方法构造最佳检验？}
\end{center}

Neyman-Pearson 定理不仅回答了"最佳检验是否存在"的问题，还给出了构造方法——它是以似然比（likelihood ratio）为基础的。

\begin{定理}[Neyman-Pearson]\label{th:neyman-pearson}
设 $X_1, \ldots, X_n$ 是来自分布 $f(x;\theta)$ 的随机样本，记似然函数为
\[
L(\theta; \mathbf{x}) = \prod_{i=1}^n f(x_i; \theta).
\]

设 $\theta'$ 和 $\theta''$ 是 $\theta$ 的两个不同取值，$\Omega = \{\theta', \theta''\}$。设 $k > 0$ 是一个常数，$C$ 是样本空间的一个子集，满足：

\begin{enumerate}
\item[(a)] $\dfrac{L(\theta';\mathbf{x})}{L(\theta'';\mathbf{x})} \leq k$ 对所有 $\mathbf{x} \in C$；
\item[(b)] $\dfrac{L(\theta';\mathbf{x})}{L(\theta'';\mathbf{x})} \geq k$ 对所有 $\mathbf{x} \in C^c$；
\item[(c)] $\alpha = P_{H_0}[\mathbf{X} \in C]$。
\end{enumerate}

则 $C$ 是检验简单假设 $H_0: \theta = \theta'$ 对 $H_1: \theta = \theta''$ 的、\textbf{大小为 $\alpha$ 的最佳临界区域}。
\end{定理}

\textbf{核心思想}：似然比 $L(\theta';\mathbf{x}) / L(\theta'';\mathbf{x})$ 衡量了数据 $\mathbf{x})$ 更支持 $H_0$ 还是更支持 $H_1$。当这个比值很小时，数据更支持 $H_1$，因此我们拒绝 $H_0$。

\textbf{物理直觉}：想象你在两个假设之间做选择。$L(\theta';\mathbf{x})$ 告诉我们"如果 $H_0$ 为真，观察到这组数据的可能性有多大"；$L(\theta'';\mathbf{x同理。如果观察到数据的可能性在 $H_0$ 下很小但在 $H_1$ 下很大，这强烈暗示 $H_1$ 才是正确的解释。}

\textbf{充分性}：值得注意的是，定理的证明表明，如果存在一个大小为 $\alpha$ 的最佳检验，那么它的形式必然是上述似然比检验。而且，这个最佳检验（如果存在）是唯一的（在连续型情况下，除了一个零概率集；在离散型中可能存在多个）。\footnote{Neyman-Pearson 引理的完整证明见附录 \cref{sec:np-lemma-proof}。}

\begin{example}[二项分布下的最佳检验]\label{ex:binomial-most-powerful}
设 $X \sim \mathrm{Binomial}(n=5, p=\theta)$，考虑检验：
\[
H_0: \theta = \tfrac{1}{2} \quad \text{versus} \quad H_1: \theta = \tfrac{3}{4}.
\]

下表给出 $f(x;1/2)$、$f(x;3/4)$ 以及它们的比值：

\begin{tabular}{c|cccccc}
$x$ & 0 & 1 & 2 & 3 & 4 & 5 \\ \hline
$f(x;1/2)$ & $1/32$ & $5/32$ & $10/32$ & $10/32$ & $5/32$ & $1/32$ \\
$f(x;3/4)$ & $1/1024$ & $15/1024$ & $90/1024$ & $270/1024$ & $405/1024$ & $243/1024$ \\
$f(x;1/2)/f(x;3/4)$ & $32/1$ & $32/3$ & $32/9$ & $32/27$ & $32/81$ & $32/243$
\end{tabular}

\textbf{问题}：构造大小 $\alpha = 1/32$ 的最佳检验。

比值 $f(x;1/2)/f(x;3/4)$ 在 $x=5$ 处取得最小值（$32/243$）。因此，最佳临界区域应包含那些使 $H_0$ 看起来极不可能而 $H_1$ 看起来极可能的数据点——即 $C = \{x = 5\}$。

验证：$P_{H_0}(X \in C) = f(5; 1/2) = 1/32 = \alpha$，
且 $P_{H_1}(X \in C) = f(5; 3/4) = 243/1024 \approx 0.237$。

如果取 $\alpha = 6/32$，则最佳临界区域为 $\{x = 4, 5\}$，因为 $x=4$ 和 $x=5$ 的比值最小。此时
\[
P_{H_1}(X \in \{4,5\}) = \frac{405}{1024} + \frac{243}{1024} = \frac{648}{1024} \approx 0.633.
\]

\textbf{关键洞察}：最佳检验包含了"最支持 $H_1$ 而最不支持 $H_0$"的数据点——这正是似然比方法告诉我们的。\footnote{二项分布最佳检验的详细分析见附录 \cref{sec:binomial-np-details}。}
\end{example}

\begin{example}[正态均值差假设的最优检验]\label{ex:normal-mean-np}
设 $X_1, \ldots, X_n \sim N(\theta, 1)$，考虑检验：
\[
H_0: \theta = 0 \quad \text{versus} \quad H_1: \theta = 1.
\]

\textbf{构造最佳检验}：似然函数为
\[
L(\theta; \mathbf{x}) = (2\pi)^{-n/2} \exp\left[-\frac{1}{2}\sum_{i=1}^n (x_i - \theta)^2\right].
\]

似然比为
\[
\frac{L(0; \mathbf{x})}{L(1; \mathbf{x})} = \exp\left(-\sum_{i=1}^n x_i + \frac{n}{2}\right).
\]

令 $\frac{L(0; \mathbf{x})}{L(1; \mathbf{x})} \leq k$，等价于
\[
\sum_{i=1}^n x_i \geq \frac{n}{2} - \log k = c.
\]

因此，最佳临界区域是 $C = \{\mathbf{x}: \sum_{i=1}^n x_i \geq c\}$，即基于样本均值 $\overline{X}$ 的检验。

\textbf{确定临界值}：当 $H_0$ 为真时，$\overline{X} \sim N(0, 1/n)$，故
\[
c_1 = \frac{c}{n} = \frac{\frac{n}{2} - \log k}{n} = \frac{1}{2} - \frac{\log k}{n}.
\]

给定显著性水平 $\alpha$，临界值为 $c_1 = q_{\alpha}(0, 1/\sqrt{n})$（$N(0,1/n)$ 的上 $\alpha$ 分位数）。

\textbf{数值示例}：设 $n = 25$，$\alpha = 0.05$，则 $c_1 = \mathbf{qnorm}(0.95, 0, 1/5) = 0.329$。

\textbf{功效计算}：当 $H_1$ 为真时，$\overline{X} \sim N(1, 1/25)$，功效为
\[
\text{Power} = P_{H_1}(\overline{X} \geq 0.329) = 1 - \Phi\left(\frac{0.329 - 1}{1/5}\right) = 1 - \Phi(-3.355) \approx 0.9996.
\]

也就是说，当 $\theta = 1$ 时，这个检验有约 $99.96\%$ 的概率正确拒绝 $H_0$。\footnote{正态均值检验的完整推导见附录 \cref{sec:normal-np-derivation}。}
\end{example}

\begin{example}[Poisson 对 Geometric 的检验]\label{ex:poisson-vs-geometric}
设 $X_1, \ldots, X_n \sim f(x)$，其中
\[
H_0: f(x) = \frac{e^{-1}}{x!}, \quad x = 0,1,2,\ldots \quad (\text{Poisson}(\lambda=1)),
\]
\[
H_1: f(x) = \left(\frac{1}{2}\right)^{x+1}, \quad x = 0,1,2,\ldots \quad (\text{Geometric}(p=1/2)).
\]

\textbf{构造最佳检验}：似然比为
\[
\frac{L(\theta'; \mathbf{x})}{L(\theta''; \mathbf{x})} = \frac{(2e^{-1})^n \cdot 2^{\sum x_i}}{\prod_{i=1}^n (x_i!)}.
\]

令 $k=1$（对应某个特定的 $\alpha$），不等式化为
\[
\frac{2^{\sum x_i}}{\prod_{i=1}^n (x_i!)} \leq \frac{e}{2}.
\]

取 $n=1$，则临界区域为 $C = \{x_1: x_1 = 0, 3, 4, 5, \ldots\}$。

\textbf{验证}：
\[
\alpha = P_{H_0}(X \in C) = 1 - \frac{e^{-1}}{1!} - \frac{e^{-1}}{2!} = 1 - \frac{2}{e} \approx 0.4482,
\]
\[
\text{Power} = P_{H_1}(X \in C) = 1 - \frac{1}{4} - \frac{1}{8} = 0.625.
\]

这符合无偏性要求（功效大于显著性水平）。\footnote{Poisson vs Geometric 检验的详细分析见附录 \cref{sec:poisson-np-details}。}
\end{example}

\textbf{无偏性（Unbiasedness）}：注意到所有最佳检验都满足一个共同性质——功效大于等于显著性水平。这意味着拒绝 $H_0$（当它为假时）的概率不小于错误拒绝它的概率。这在经济决策中是完全合理的——如果检验的功效低于 $\alpha$，那它的行为还不如随机猜测。

\begin{definition}[无偏检验]\label{def:unbiased-test}
设 $C$ 是显著性水平为 $\alpha$ 的检验的临界区域。如果对所有 $\theta \in \omega_1$，
\[
P_\theta(\mathbf{X} \in C) \geq \alpha,
\]
则称该检验是\textbf{无偏的（unbiased）}。
\end{definition}

推论 8.1.1 表明 Neyman-Pearson 最佳检验是无偏的——在 $H_0$ 为假时拒绝它的概率不低于在 $H_0$ 为真时错误拒绝它的概率。\footnote{推论 8.1.1 的证明及无偏性的讨论见附录 \cref{sec:unbiased-np-corollary}。}

\section{8.2 Uniformly Most Powerful Tests}\label{sec:8.2}

\subsection{从简单假设到复合假设}

上一节的 Neyman-Pearson 定理解决的是"简单假设对简单假设"的问题——$H_0$ 和 $H_1$ 都各只有一个参数值。但在实际中，备择假设通常是复合的（composite），例如：

\[
H_0: \theta \leq \theta' \quad \text{versus} \quad H_1: \theta > \theta'.
\]

这里 $H_1$ 包含无穷多个参数值。对于这样的问题，我们自然希望找到一个检验，它在 $H_0$ 的每一个边界点上都是最佳的。

\begin{center}
\textit{能否找到一个检验，它对 $\theta > \theta'$ 中的每一个 $\theta$ 都是最佳检验？}
\end{center}

这就是"一致最大功效（Uniformly Most Powerful，UMP）"检验的思想。

\begin{definition}[UMP 检验]\label{def:ump-test}
设 $C$ 是显著性水平为 $\alpha$ 的检验的临界区域。如果 $C$ 是 $H_0$ 对 $H_1$ 中每一个简单假设的最佳临界区域，则称 $C$ 为\textbf{大小为 $\alpha$ 的一致最大功效临界区域（UMP critical region）}，相应的检验称为\textbf{UMP 检验（Uniformly Most Powerful test）}。
\end{definition}

但 UMP 检验并不总是存在。例 8.2.3 展示了：对于正态均值 $\sigma^2 = 1$ 已知时，检验 $H_0: \theta = \theta'$ 对 $H_1: \theta \neq \theta'$（双向备择）不存在 UMP 检验——因为对 $\theta'' > \theta'$ 的最佳临界区域（拒绝域在右尾）与对 $\theta'' < \theta'$ 的最佳临界区域（拒绝域在左尾）形式完全不同，无法统一。

这个例子揭示了一个深刻的事实：\textbf{UMP 检验只在单向备择假设（one-sided alternatives）下才可能存在}。

\subsection{单调似然比与 Karlin-Rubin 定理}

那么，在什么条件下 UMP 检验存在呢？答案藏在一种叫做"单调似然比（Monotone Likelihood Ratio，MLR）"的性质中。

\begin{definition}[单调似然比]\label{def:mlr}
设 $L(\theta, \mathbf{x})$ 是似然函数。如果对 $\theta_1 < \theta_2$，比值
\[
\frac{L(\theta_1; \mathbf{x})}{L(\theta_2; \mathbf{x})}
\]
是统计量 $y = u(\mathbf{x})$ 的单调递减函数，则称似然函数在 $y = u(\mathbf{x})$ 中有\textbf{单调似然比（monotone likelihood ratio, mlr）}。
\end{definition}

\textbf{为什么 MLR 如此重要？} MLR 性质确保了对较大的 $y$（即更支持备择假设的数据），我们统一拒绝 $H_0$——而不需要针对每一个具体的 $\theta''$ 分别构造最佳检验。

\begin{定理}[Karlin-Rubin 定理]\label{th:karlin-rubin}
设随机样本来自密度 $f(x;\theta)$，似然函数 $L(\theta, \mathbf{x})$ 在统计量 $Y = u(\mathbf{X})$ 中有单调似然比。考虑单向假设：
\[
H_0: \theta \leq \theta' \quad \text{versus} \quad H_1: \theta > \theta'.
\]

则对任意显著性水平 $\alpha$，UMP 检验为：
\[
\text{Reject } H_0 \quad \text{if} \quad Y \geq c_Y,
\]
其中临界值 $c_Y$ 满足 $P_{\theta'}[Y \geq c_Y] = \alpha$。
\end{定理}

\textbf{核心洞察}：由于 MLR 性质，对任何 $\theta'' > \theta'$，最佳临界区域的形式都是 $Y \geq c$——它们恰好是同一个区域！因此这个区域对所有 $\theta'' > \theta'$ 都是最佳的，从而是一致最大功效的。

\textbf{证明思路}：对固定的 $\theta'' > \theta'$，由 Neyman-Pearson 定理，最佳临界区域由似然比 $L(\theta';\mathbf{x}) / L(\theta'';\mathbf{x}) \leq k$ 给出。由于 MLR 性质，这个不等式等价于 $Y \geq g^{-1}(k)$——即只依赖于数据的形状，而不依赖于具体的 $\theta''$ 值。\footnote{Karlin-Rubin 定理的完整证明见附录 \cref{sec:karlin-rubin-proof}。}

\begin{example}[正态方差 UMP 检验]\label{ex:normal-variance-ump}
设 $X_1, \ldots, X_n \sim N(0, \theta)$，其中 $\theta > 0$ 是方差。考虑检验：
\[
H_0: \theta = \theta' \quad \text{versus} \quad H_1: \theta > \theta'.
\]

似然函数为
\[
L(\theta; \mathbf{x}) = \left(\frac{1}{2\pi\theta}\right)^{n/2} \exp\left\{-\frac{1}{2\theta}\sum_{i=1}^n x_i^2\right\}.
\]

似然比为
\[
\frac{L(\theta'; \mathbf{x})}{L(\theta''; \mathbf{x})} = \left(\frac{\theta''}{\theta'}\right)^{n/2} \exp\left[-\left(\frac{\theta'' - \theta'}{2\theta'\theta''}\right)\sum_{i=1}^n x_i^2\right].
\]

对 $\theta'' > \theta'$，这是 $\sum x_i^2$ 的递减函数，因此似然函数在 $Y = \sum X_i^2$ 中有 MLR。

由 Karlin-Rubin 定理，UMP 检验为：
\[
\text{Reject } H_0 \quad \text{if} \quad \sum_{i=1}^n X_i^2 \geq c,
\]
其中 $c$ 满足 $P_{\theta'}[\sum X_i^2 \geq c] = \alpha$。

\textbf{数值示例}：设 $n = 15$，$\alpha = 0.05$，$\theta' = 3$。当 $H_0$ 为真时，$\sum X_i^2 / \theta' \sim \chi^2(n)$，故
\[
\frac{c}{3} = \chi^2_{0.95}(15) = 24.996 \quad \Rightarrow \quad c = 74.988.
\]

即当 $\sum x_i^2 \geq 74.988$ 时拒绝 $H_0: \theta = 3$。\footnote{正态方差 UMP 检验的详细推导见附录 \cref{sec:normal-variance-ump-derivation}。}
\end{example}

\begin{example}[Bernoulli 分布的 UMP 检验]\label{ex:bernoulli-ump}
设 $X_1, \ldots, X_n \sim \mathrm{Bernoulli}(\theta)$，考虑检验：
\[
H_0: \theta \leq \theta' \quad \text{versus} \quad H_1: \theta > \theta'.
\]

似然比为
\[
\frac{L(\theta'; \mathbf{x})}{L(\theta''; \mathbf{x})} = \left[\frac{\theta'(1-\theta'')}{\theta''(1-\theta')}\right]^{\sum x_i} \left(\frac{1-\theta'}{1-\theta''}\right)^n.
\]

由于 $\theta' < \theta''$ 时 $\theta'(1-\theta'')/\theta''(1-\theta') < 1$，该比值是 $Y = \sum X_i$ 的递减函数。因此似然函数在 $Y = \sum X_i$ 中有 MLR，UMP 检验为：
\[
\text{Reject } H_0 \quad \text{if} \quad Y = \sum_{i=1}^n X_i \geq c,
\]
其中 $c$ 由 $P_{\theta'}[Y \geq c] = \alpha$ 确定。

这正是我们直观上期望的：观察到更多的成功次数，就越有理由拒绝"成功率不高"的原假设。\footnote{Bernoulli UMP 检验的 MLR 证明见附录 \cref{sec:bernoulli-ump-derivation}。}
\end{example}

\textbf{正规指数族（Regular Exponential Family）}：其实 Bernoulli 分布、正态分布（方差已知）、Poisson 分布都属于一个更广泛的族——正规指数族。在这个族中，只要自然参数是单调的，就存在 UMP 检验。具体地，如果分布有形式
\[
f(x;\theta) = \exp[p(\theta)K(x) + H(x) + q(\theta)],
\]
且 $p(\theta)$ 是 $\theta$ 的递增函数，则对假设 $H_0: \theta \leq \theta'$ 对 $H_1: \theta > \theta'$，UMP 检验存在，基于统计量 $Y = \sum_{i=1}^n K(X_i)$。\footnote{正规指数族的 UMP 检验存在性讨论见附录 \cref{sec:exp-family-ump}。}

\section{8.3 Likelihood Ratio Tests}\label{sec:8.3}

\subsection{似然比检验的动机}

在前两节中，我们学会了如何在简单假设对简单假设时构造最佳检验，以及在具有 MLR 性质的族中对单向复合假设构造 UMP 检验。但现实中很多重要问题超出了这些范畴：

\begin{itemize}
\item 检验正态均值 $H_0: \mu = \mu_0$（$\sigma^2$ 未知）对 $H_1: \mu \neq \mu_0$——不存在 UMP 检验（双向备择）；
\item 检验两正态均值是否相等 $H_0: \mu_1 = \mu_2$——涉及多于一个参数；
\item 一般复合假设——没有 MLR 性质可供利用。
\end{itemize}

在这些情况下，我们转向一种通用但不一定最优的构造检验的方法——\textbf{似然比检验（Likelihood Ratio Test，LRT）}。

\textbf{为什么即使 LRT 不是最优的，我们仍然使用它？} 历史上，Neyman 和 Pearson 引入 LRT 的动机正是为了突破最佳检验只存在于简单假设的限制。LRT 继承了 Neyman-Pearson 定理的似然比思想（用似然比衡量数据对假设的支持程度），但将比较扩展到了复合假设。虽然在一般情况下 LRT 不能保证最优性，但它是渐近最优的——在样本量趋向无穷时，LRT 具有许多优良性质。\footnote{LRT 与 Neyman-Pearson 引理的关系讨论见附录 \cref{sec:lrt-np-relationship}。}

\subsection{LRT 的定义}

设 $X$ 的密度为 $f(\mathbf{x}; \boldsymbol{\theta})$，参数 $\boldsymbol{\theta} \in \Omega \subset \mathbb{R}^p$。考虑复合假设：
\[
H_0: \boldsymbol{\theta} \in \omega \quad \text{versus} \quad H_1: \boldsymbol{\theta} \in \Omega \cap \omega^c, \tag{8.3.1}\label{eq:general-hypotheses}

\]
其中 $\omega \subset \Omega$ 是参数空间的某个子集。

\begin{definition}[似然比]\label{def:likelihood-ratio}
\textbf{似然比（Likelihood Ratio）}定义为：
\[
\Lambda(\mathbf{x}) = \frac{\sup_{\omega} L(\boldsymbol{\theta}; \mathbf{x})}{\sup_{\Omega} L(\boldsymbol{\theta}; \mathbf{x})} = \frac{L(\hat{\boldsymbol{\omega}}; \mathbf{x})}{L(\hat{\boldsymbol{\Omega}}; \mathbf{x})},
\]
其中 $\hat{\boldsymbol{\omega}}$ 是在约束 $\boldsymbol{\theta} \in \omega$ 下的 MLE，$\hat{\boldsymbol{\Omega}}$ 是在无约束 $\boldsymbol{\theta} \in \Omega$ 下的 MLE。
\end{definition}

\textbf{直观含义}：分子衡量数据在 $H_0$ 约束下能达到的最大似然，分母衡量数据在无约束情况下能达到的最大似然。如果 $H_0$ 为真，这两个值应该差不多，比值接近 1；如果 $H_0$ 为假，无约束的似然会显著大于约束下的似然，比值会变小。因此我们\textbf{在小比值时拒绝 $H_0$}：
\[
\text{Reject } H_0 \quad \text{if} \quad \Lambda(\mathbf{x}) \leq \lambda_0 < 1.
\]

\textbf{LRT 的一个关键性质}：在 $H_0$ 下，$-2\log\Lambda$ 渐近服从 $\chi^2$ 分布，自由度为 $\dim(\Omega) - \dim(\omega)$（在一定正则条件下）。这使得我们可以近似地确定临界值。

\textbf{LRT 与 NP 引理的关系}：当 $\Omega = \{\theta', \theta''\}$（只有两点）时，LRT 还原为 NP 似然比检验，此时它是最优的。但一般来说，LRT 不一定是最优的。

接下来我们把 LRT 应用到几个具体的重要情形——正态均值差和正态方差检验。

\subsection{8.3.1 正态均值差的 LRT}\label{sec:8.3.1}

\begin{example}[两样本 $t$ 检验]\label{ex:two-sample-t-test}
设 $X_1, \ldots, X_n \sim N(\theta_1, \theta_3)$ 和 $Y_1, \ldots, Y_m \sim N(\theta_2, \theta_3)$，两组样本独立，共同方差 $\theta_3$ 未知。考虑检验：
\[
H_0: \theta_1 = \theta_2 \quad \text{versus} \quad H_1: \theta_1 \neq \theta_2.
\]

\textbf{构造 LRT}：在 $H_0$ 下（$\theta_1 = \theta_2 = \mu$），参数的 MLE 为
\[
\hat{\mu} = \frac{n\bar{x} + m\bar{y}}{n+m}, \quad \hat{\sigma}_0^2 = \frac{1}{n+m}\left[\sum_{i=1}^n (x_i - \hat{\mu})^2 + \sum_{i=1}^m (y_i - \hat{\mu})^2\right].
\]

在无约束下（$H_1$），MLE 分别为 $\bar{x}$、$\bar{y}$ 和
\[
\hat{\sigma}^2 = \frac{1}{n+m}\left[\sum_{i=1}^n (x_i - \bar{x})^2 + \sum_{i=1}^m (y_i - \bar{y})^2\right].
\]

似然比为
\[
\Lambda = \frac{L(\hat{\omega})}{L(\hat{\Omega})} = \left(\frac{\hat{\sigma}^2}{\hat{\sigma}_0^2}\right)^{(n+m)/2}.
\]

\textbf{与 $t$ 统计量的联系}：可以证明
\[
\Lambda^{2/(n+m)} = \frac{(n+m-2)}{(n+m-2) + T^2},
\]
其中
\[
T = \sqrt{\frac{nm}{n+m}} \frac{\bar{x} - \bar{y}}{S_p}, \quad S_p^2 = \frac{(n-1)S_x^2 + (m-1)S_y^2}{n+m-2}.
\]

$T$ 服从自由度为 $n+m-2$ 的 $t$ 分布（在 $H_0$ 下）。由于 $\Lambda \leq \lambda_0 \Leftrightarrow |T| \geq c$，LRT 就是经典的\textbf{两样本 $t$ 检验}。

\textbf{数值示例}：设 $n = 10$，$m = 6$，$\alpha = 0.05$，则临界值
\[
c = t_{0.975}(14) = 2.1448.
\]

如果 $|t| \geq 2.1448$，拒绝 $H_0$（认为两个总体均值不相等）。\footnote{两样本 $t$ 检验的完整推导见附录 \cref{sec:two-sample-t-derivation}。}
\end{example}

\textbf{功效函数}：对于非中心 $t$ 分布，当 $\theta_1 \neq \theta_2$ 时，$T$ 服从非中心 $t$ 分布，非中心参数为
\[
\delta = \sqrt{\frac{nm}{n+m}} \frac{\theta_1 - \theta_2}{\sigma}.
\]

给定 $\theta_3 = 100$、$n = 20$、$m = 15$、$\alpha = 0.05$、$\Delta = 5 = \theta_1 - \theta_2$，临界值 $t_{0.975}(33) = 2.0345$，非中心参数 $\delta_2 = 1.4639$，功效约为
\[
1 - \text{pt}(2.0345, 33, \text{ncp}=1.4639) + \text{pt}(-2.0345, 33, \text{ncp}=1.4639) \approx 0.2954.
\]

即约 $29.5\%$ 的概率检测出均值为 5 的差异。\footnote{非中心 $t$ 分布的功效计算见附录 \cref{sec:noncentral-t-power}。}

\textbf{稳健性注记}：$t$ 检验在正态性假设不满足时表现如何？在大样本下，由 Slutsky 定理，即使底层分布非正态，只要方差有限，$t$ 统计量仍然渐近正态——因此 $t$ 检验是\textbf{渐近正确的}。然而，在有限样本且严重非正态（特别是高度偏态）时，$t$ 检验可能失效——实际显著性水平会偏离名义水平。\footnote{非正态情况下 $t$ 检验的稳健性讨论见附录 \cref{sec:t-test-robustness}。}

\subsection{8.3.2 正态方差的 LRT}\label{sec:8.3.2}

\begin{example}[两方差相等的 F 检验]\label{ex:f-test-variance}
设 $X_1, \ldots, X_n \sim N(\theta_1, \theta_3)$ 和 $Y_1, \ldots, Y_m \sim N(\theta_2, \theta_4)$，独立样本，考虑检验：
\[
H_0: \theta_3 = \theta_4 \quad \text{versus} \quad H_1: \theta_3 \neq \theta_4.
\]

可以证明（练习 8.3.11），LRT 统计量可以化为
\[
F = \frac{\sum_{i=1}^n (X_i - \bar{X})^2 / (n-1)}{\sum_{i=1}^m (Y_i - \bar{Y})^2 / (m-1)}. \tag{8.3.2}\label{eq:f-statistic-variance}

\]

在 $H_0$ 下，$F \sim F(n-1, m-1)$。因此双向 LRT 为：
\[
\text{Reject } H_0 \quad \text{if} \quad F \leq c_1 \quad \text{or} \quad F \geq c_2,
\]
其中 $c_1$ 和 $c_2$ 满足 $P(F \leq c_1) = P(F \geq c_2) = \alpha_1/2$。

\textbf{重要警告}：与 $t$ 检验不同，$F$ 检验对正态性假设是\textbf{敏感}的。在非正态情况下，$F$ 检验的名义水平和实际水平可能相差甚远（可高达 $0.80$），因此不应在正态性未经检验的情况下使用 $F$ 检验。\footnote{方差 $F$ 检验的稳健性问题见附录 \cref{sec:f-test-robustness}。}
\end{example}

\section{8.4 *Sequential Probability Ratio Test}\label{sec:8.4}

\subsection{固定样本量检验的局限性}

到目前为止，我们讨论的所有检验都有一个共同假设：样本量 $n$ 是预先固定的。但在实际应用中，固定样本量检验存在一个根本性的低效：

\begin{center}
\textit{为什么要等到收集完所有 $n$ 个样本才做决策？为什么不能边收集边决策？}
\end{center}

考虑一个质量控制场景：你需要判断一条生产线是否正常（$\theta = \theta_0$）。固定样本量方法要求你先收集 $n$ 个样本，然后做决策。但如果在第 5 个样本时，数据就已经极其强烈地支持 $H_1$（生产线已严重偏离正常状态），继续等待其余 $n-5$ 个样本就是在浪费时间和资源。

1930 年代，Abraham Wald 提出了一个革命性的思想：\textbf{允许样本量根据数据动态确定}——这就是序贯分析（sequential analysis）的诞生。

Wald 的序贯概率比检验（SPRT）不仅在概念上优雅，而且在功效和样本量之间实现了最优权衡。

\subsection{SPRT 的构造}

设 $X_1, X_2, \ldots$ 是 iid 随机变量序列，考虑检验：
\[
H_0: \theta = \theta' \quad \text{versus} \quad H_1: \theta = \theta'',
\]
其中 $\theta'$ 和 $\theta''$ 是给定的具体值。

\textbf{核心思想}：在每个步骤 $n$，计算似然比
\[
\Lambda_n = \frac{L(\theta'; x_1, \ldots, x_n)}{L(\theta''; x_1, \ldots, x_n)} = \prod_{i=1}^n \frac{f(x_i; \theta')}{f(x_i; \theta'')}.

\]

如果 $\Lambda_n$ 非常小（即数据强烈支持 $H_1$），我们就提前停止并拒绝 $H_0$；如果 $\Lambda_n$ 非常大（即数据强烈支持 $H_0$），我们就提前停止并接受 $H_0$；如果 $\Lambda_n$ 处于中间位置，我们就继续观察下一个样本。

具体地，选取两个常数 $A < B$（通常 $A = \beta/(1-\alpha)$，$B = (1-\beta)/\alpha$），SPRT 的决策规则为：

\[
\text{在第 $n$ 步，停止采样并：}
\begin{cases}
\text{拒绝 $H_0$（接受 $H_1$）} & \text{if } \Lambda_n \leq A, \\
\text{接受 $H_0$（拒绝 $H_1$）} & \text{if } \Lambda_n \geq B, \\
\text{继续采样} & \text{if } A < \Lambda_n < B.
\end{cases}
\tag{8.4.1}\label{eq:sprt-rule}

\]

\textbf{Wald 的关键近似}：在实际应用中，Wald 证明了可以通过近似 $A \approx \beta/\alpha$ 和 $B \approx (1-\beta)/(1-\alpha)$ 来设置常数，其中 $\alpha$ 和 $\beta$ 分别是目标第一类和第二类错误概率。

\begin{example}[Bernoulli 的 SPRT]\label{ex:bernoulli-sprt}
设 $X_i \sim \mathrm{Bernoulli}(\theta)$，检验 $H_0: \theta = 1/3$ 对 $H_1: \theta = 2/3$。

似然比为
\[
\Lambda_n = \frac{(1/3)^{\sum x_i} (2/3)^{n-\sum x_i}}{(2/3)^{\sum x_i} (1/3)^{n-\sum x_i}} = 2^{n - 2\sum x_i}.
\]

取对数并整理，不等式 $A < \Lambda_n < B$ 等价于
\[
c_0(n) = \frac{n}{2} - \frac{\log_2 B}{2} < \sum_{i=1}^n x_i < \frac{n}{2} - \frac{\log_2 A}{2} = c_1(n).
\]

\textbf{决策规则}：当 $\sum x_i \geq c_1(n)$ 时，拒绝 $H_0$；当 $\sum x_i \leq c_0(n)$ 时，接受 $H_0$；否则继续采样。

\textbf{序贯抽样的图示}：想象一条随机游走——每观察一次 $X_i = 1$，我们就向右跳 $1$ 个单位；每观察一次 $X_i = 0$，我们就向左跳 $1$ 个单位。只要这条路径处于上下边界之间，我们就继续。一旦路径碰到上界（强烈支持 $H_1$）或下界（强烈支持 $H_0$），我们就停止并做出相应决策。\footnote{Bernoulli SPRT 的详细分析见附录 \cref{sec:bernoulli-sprt-details}。}
\end{example}

\begin{example}[正态均值的 SPRT]\label{ex:normal-sprt}
设 $X_1, X_2, \ldots \sim N(\theta, 100)$，检验 $H_0: \theta = 75$ 对 $H_1: \theta = 78$，目标是使 $\alpha \approx \beta \approx 0.10$。

取 $A = 1/9$、$B = 9$。似然比为
\[
\Lambda_n = \exp\left(-\frac{6\sum x_i - 459n}{200}\right).
\]

不等式 $A < \Lambda_n < B$ 化为
\[
c_0(n) = \frac{153}{2}n - \frac{100}{3}\log 9 < \sum_{i=1}^n x_i < \frac{153}{2}n + \frac{100}{3}\log 9 = c_1(n).
\]

决策规则：在第 $n$ 步，当 $\sum x_i \geq c_1(n)$ 时拒绝 $H_0$（认为 $\theta = 78$）；当 $\sum x_i \leq c_0(n)$ 时接受 $H_0$（认为 $\theta = 75$）；否则继续采样。\footnote{正态 SPRT 的完整推导见附录 \cref{sec:normal-sprt-details}。}
\end{example}

\textbf{SPRT 的吸引力}：为什么 SPRT 如此重要？因为它在固定两类错误概率的同时，通常使用比固定样本量检验更少的样本。Wald 的近似理论表明，当 $H_0$ 或 $H_1$ 为真时，SPRT 的平均样本量（Average Sample Number，ASN）比相应的固定样本量检验小得多——有时能节省 $50\%$ 以上的样本。

\textbf{SPRT 与质量控制}：SPRT 的思想在工业质量控制中有广泛应用。CUSUM（累积和）控制图就是 SPRT 的一个变体——它持续累积标准化后的偏离量 $\sum (Z_i - 0.5)$（其中 $Z_i$ 是标准化后的检验统计量），一旦超出控制限就发出警报。与固定间隔的 Shewhart 控制图相比，CUSUM 能更快检测到过程均值的微小漂移。

\section{8.5 *Minimax and Classification Procedures}\label{sec:8.5}

\subsection{8.5.1 极小化极大检验}

在第 7 章的决策论框架中，我们介绍了极小化极大（minimax）原则：将最大风险最小化。将这个思想应用到假设检验中，会得到一个有趣的结论——\textbf{Neyman-Pearson 最佳检验在某种意义上也是极小化极大的}。

考虑简单假设 $H_0: \theta = \theta'$ 对 $H_1: \theta = \theta''$，损失函数满足 $\mathcal{L}(\theta', \theta') = \mathcal{L}(\theta'', \theta'') = 0$（正确决策无损失），而 $\mathcal{L}(\theta', \theta'') > 0$ 和 $\mathcal{L}(\theta'', \theta') > 0$（错误决策有损失）。\footnote{极小化极大检验的决策论框架见附录 \cref{sec:minimax-test-framework}。}

此时风险函数为：
\[
R(\theta', C) = \mathcal{L}(\theta', \theta'') \cdot \alpha, \quad R(\theta'', C) = \mathcal{L}(\theta'', \theta') \cdot \beta,

\]
其中 $\alpha$ 是显著性水平，$\beta$ 是 Type II 错误概率。

极小化极大原则要求选择 $C$ 使得 $\max\{R(\theta', C), R(\theta'', C)\}$ 最小。可以证明，当 $k$ 满足
\[
\mathcal{L}(\theta', \theta'') \int_C L(\theta') = \mathcal{L}(\theta'', \theta') \int_{C^c} L(\theta'')

\]
时（即两类错误造成的风险相等），临界区域 $C$ 是极小化极大的。这个条件与 NP 似然比条件是一致的——本质上，最佳检验同时也是极小化极大检验（当两类错误的相对代价被适当选择时）。\footnote{极小化极大检验的完整证明见附录 \cref{sec:minimax-test-proof}。}

\begin{example}[正态均值的极小化极大检验]\label{ex:minimax-normal}
设 $X_1, \ldots, X_{100} \sim N(\theta, 100)$，检验 $H_0: \theta = 75$ 对 $H_1: \theta = 78$，损失为 $\mathcal{L}(75, 78) = 3$、$\mathcal{L}(78, 75) = 1$。

临界区域形如 $\bar{x} \geq c$，其中 $c$ 满足：
\[
3 \cdot P_{\theta=75}(\bar{X} \geq c) = P_{\theta=78}(\bar{X} < c).

\]
由于 $\bar{X} \sim N(\theta, 1)$，这化为
\[
3[1 - \Phi(c-75)] = \Phi(c-78).

\]

求解得 $c \approx 76.8$。此时显著性水平 $\alpha = 1 - \Phi(1.8) \approx 0.036$，功效为 $1 - \Phi(-1.2) \approx 0.885$。\footnote{极小化极大检验的数值求解（Newton 法）见附录 \cref{sec:minimax-newton}。}
\end{example}

\subsection{8.5.2 分类问题}

极小化极大检验理论在分类（classification）问题中有直接应用。考虑一个二分类问题：我们观察到 $(X, Y)$，需要判断它来自分布 $f(x,y;\theta')$ 还是 $f(x,y;\theta'')$。

这本质上就是假设检验问题——如果接受 $H_0$，就把 $(X,Y)$ 分类为来自 $\theta'$ 类；如果拒绝 $H_0$，就分类为来自 $\theta''$ 类。

由 NP 定理，最佳分类规则是似然比规则：
\[
\frac{f(x,y;\theta')}{f(x,y;\theta'')} \leq k \quad \Longrightarrow \quad \text{分类为 $\theta''$ 类}.

\]

若两类的先验概率为 $\pi'$ 和 $\pi''$，则最优规则取 $k = \pi''/\pi'$（当两类等可能时 $k=1$）。

\textbf{Fisher 线性判别函数}：在正态分布的特殊情况下，似然比不等式化为一个线性函数 $ax + by \leq c$——这就是著名的 Fisher 线性判别函数。当参数未知时，用训练样本的均值和方差代替，就得到了 Fisher 判别函数。\footnote{Fisher 线性判别函数的完整推导见附录 \cref{sec:fisher-discriminant}。}

\begin{Example}[二元正态的线性判别]\label{ex:bivariate-normal-classification}
设 $(X,Y)$ 来自二元正态分布，参数为 $(\mu_1', \mu_2', \sigma_1^2, \sigma_2^2, \rho')$ 或 $(\mu_1'', \mu_2'', \sigma_1^2, \sigma_2^2, \rho'')$。则似然比不等式化为
\[
ax + by \leq c,

\]
其中 $a$ 和 $b$ 是由两个总体的均值向量和协方差矩阵决定的常数。这给出了一个\textbf{线性决策边界}——在二维空间中是一条直线，将平面分为两部分，分别对应两类。\footnote{二元正态线性判别的详细推导见附录 \cref{sec:bivariate-discriminant}。}
\end{Example}

\section*{附录：公式推导}\label{sec:appendix-ch8}

本附录包含 Chapter 8 中省略的完整推导。

\subsection{Neyman-Pearson 引理的证明}\label{sec:np-lemma-proof}

\textbf{背景（Background）}：我们需要证明，当临界区域 $C$ 满足 NP 条件（a）-(c) 时，它是在固定显著性水平 $\alpha$ 下功效最高的检验。

\textbf{已知条件（Given）}：
\begin{itemize}
\item $L(\theta'; \mathbf{x}) = \prod_{i=1}^n f(x_i; \theta')$ 是 $H_0$ 下的似然函数
\item $L(\theta''; \mathbf{x}) = \prod_{i=1}^n f(x_i; \theta'')$ 是 $H_1$ 下的似然函数
\item 临界区域 $C$ 满足：
\[
\frac{L(\theta';\mathbf{x})}{L(\theta'';\mathbf{x})} \leq k, \quad \forall \mathbf{x} \in C;
\quad
\frac{L(\theta';\mathbf{x})}{L(\theta'';\mathbf{x})} \geq k, \quad \forall \mathbf{x} \in C^c.
\]
\item 任意其他满足 $P_{\theta'}[\mathbf{X} \in A] = \alpha$ 的区域 $A$。
\end{itemize}

\textbf{目标（Goal）}：证明 $P_{\theta''}[\mathbf{X} \in C] \geq P_{\theta''}[\mathbf{X} \in A]$。

\textbf{推导步骤（Derivation Steps）}：

1. 写出功效差：
\[
\int_C L(\theta'') - \int_A L(\theta'') = \int_{C \cap A^c} L(\theta'') - \int_{A \cap C^c} L(\theta'').

\tag{8.A.1}\label{eq:np-proof-eq1}

\]

2. 由条件 (a)：在 $C$ 上 $L(\theta'') \geq \frac{1}{k}L(\theta')$，故
\[
\int_{C \cap A^c} L(\theta'') \geq \frac{1}{k} \int_{C \cap A^c} L(\theta').

\tag{8.A.2}\label{eq:np-proof-eq2}

\]

3. 由条件 (b)：在 $C^c$ 上 $L(\theta'') \leq \frac{1}{k}L(\theta')$，故
\[
\int_{A \cap C^c} L(\theta'') \leq \frac{1}{k} \int_{A \cap C^c} L(\theta').

\tag{8.A.3}\label{eq:np-proof-eq3}

\]

4. 将 (8.A.2) 和 (8.A.3) 代入 (8.A.1)：
\[
\int_C L(\theta'') - \int_A L(\theta'') \geq \frac{1}{k}\left[\int_{C \cap A^c} L(\theta') - \int_{A \cap C^c} L(\theta')\right].

\tag{8.A.4}\label{eq:np-proof-eq4}

\]

5. 计算括号内项：
\[
\int_{C \cap A^c} L(\theta') - \int_{A \cap C^c} L(\theta') = \int_C L(\theta') - \int_A L(\theta') = \alpha - \alpha = 0.

\tag{8.A.5}\label{eq:np-proof-eq5}

\]

6. 由 (8.A.4) 和 (8.A.5) 得 $\int_C L(\theta'') - \int_A L(\theta'') \geq 0$，即
\[
P_{\theta''}[\mathbf{X} \in C] \geq P_{\theta''}[\mathbf{X} \in A].

\]

证毕。

\subsection{正态均值 NP 检验的完整推导}\label{sec:normal-np-derivation}

\textbf{背景（Background）}：在例 8.1.2 中，我们构造了正态均值差假设的 NP 检验。本节补充完整推导。

\textbf{目标（Goal）}：对 $X_1, \ldots, X_n \sim N(\theta, 1)$，检验 $H_0: \theta = 0$ 对 $H_1: \theta = 1$，求大小为 $\alpha$ 的最佳检验。

\textbf{推导步骤（Derivation Steps）}：

1. 写出似然函数：
\[
L(\theta; \mathbf{x}) = (2\pi)^{-n/2} \exp\left[-\frac{1}{2}\sum_{i=1}^n (x_i - \theta)^2\right].

\]

2. 似然比：
\[
\frac{L(0; \mathbf{x})}{L(1; \mathbf{x})} = \frac{\exp\left[-\frac{1}{2}\sum x_i^2\right]}{\exp\left[-\frac{1}{2}\sum (x_i-1)^2\right]} = \exp\left[-\sum_{i=1}^n x_i + \frac{n}{2}\right].

\]

3. 令 $\frac{L(0; \mathbf{x})}{L(1; \mathbf{x})} \leq k$：
\[
\exp\left[-\sum x_i + \frac{n}{2}\right] \leq k \quad \Longleftrightarrow \quad \sum_{i=1}^n x_i \geq \frac{n}{2} - \log k = c.

\]

4. 当 $H_0$ 为真时，$\overline{X} \sim N(0, 1/n)$，故临界值
\[
c_1 = \frac{c}{n} = \frac{1}{2} - \frac{\log k}{n} = q_{\alpha}(0, 1/\sqrt{n}).

\]

5. 功效函数：当 $\theta = 1$ 时，$\overline{X} \sim N(1, 1/n)$，
\[
\text{Power}(\theta=1) = P_{\theta=1}\left(\overline{X} \geq q_{\alpha}(0, 1/\sqrt{n})\right) = 1 - \Phi\left(\frac{q_{\alpha}(0, 1/\sqrt{n}) - 1}{1/\sqrt{n}}\right).

\]

证毕。

\subsection{Karlin-Rubin 定理的证明}\label{sec:karlin-rubin-proof}

\textbf{背景（Background）}：证明当似然函数有 MLR 性质时，单向复合假设存在 UMP 检验。

\textbf{已知条件（Given）}：
\begin{itemize}
\item 似然函数 $L(\theta, \mathbf{x})$ 在 $Y = u(\mathbf{x})$ 中有 MLR：对 $\theta_1 < \theta_2$，$L(\theta_1;\mathbf{x})/L(\theta_2;\mathbf{x}) = g(Y)$，其中 $g$ 是递减函数。
\item 单向假设：$H_0: \theta \leq \theta'$ 对 $H_1: \theta > \theta'$。
\end{itemize}

\textbf{目标（Goal）}：证明检验 $\phi(\mathbf{x}) = 1$ 当 $Y \geq c_Y$，其中 $c_Y$ 满足 $P_{\theta'}[Y \geq c_Y] = \alpha$，是 UMP level $\alpha$ 检验。

\textbf{推导步骤（Derivation Steps）}：

1. \textbf{对简单假设 $H_0': \theta = \theta'$ 对 $\theta'' > \theta'$ 的最优性}：由 NP 定理，最佳临界区域由
\[
\frac{L(\theta';\mathbf{x})}{L(\theta'';\mathbf{x})} \leq k

\]
给出。由于 MLR 性质，$g(Y) \leq k \Longleftrightarrow Y \geq g^{-1}(k)$。取 $c_Y = g^{-1}(k)$，则对任意 $\theta'' > \theta'$，临界区域 $C = \{Y \geq c_Y\}$ 是 $H_0': \theta = \theta'$ 对 $\theta''$ 的最佳检验。

2. \textbf{对所有 $\theta'' > \theta'$ 的一致性}：由于 NP 最佳临界区域的形式只依赖于 $Y \geq c_Y$（而不依赖于具体的 $\theta''$），$C = \{Y \geq c_Y\}$ 对所有 $\theta'' > \theta'$ 都是最佳的。

3. \textbf{无偏性（保证 level 不超过 $\alpha$）}：设 $\gamma_Y(\theta) = P_\theta(Y \geq c_Y)$ 为功效函数。对 $\theta_1 < \theta_2$，由第 1 步，$C$ 是 $\theta_1$ 对 $\theta_2$ 的最佳检验，故 $\gamma_Y(\theta_2) \geq \gamma_Y(\theta_1)$。因此 $\gamma_Y(\theta)$ 是 $\theta$ 的非降函数，$\max_{\theta \leq \theta'} \gamma_Y(\theta) = \gamma_Y(\theta') = \alpha$。故该检验是 level $\alpha$ 的。

综上，$C = \{Y \geq c_Y\}$ 是UMP level $\alpha$ 检验。证毕。

\section*{Exercises}\label{sec:exercises-ch8}

\begin{exercise}{8.1.4 — Best Critical Region for Normal Variance}
Let $X_1, X_2, \ldots, X_{10}$ be a random sample of size 10 from a normal distribution $N(0, \sigma^2)$. Find a best critical region of size $\alpha = 0.05$ for testing $H_0: \sigma^2 = 1$ against $H_1: \sigma^2 = 2$. Is this a best critical region of size 0.05 for testing $H_0: \sigma^2 = 1$ against $H_1: \sigma^2 = 4$? Against $H_1: \sigma^2 = \sigma_1^2 > 1$?
\end{exercise}

\begin{exercise}{8.2.7 — UMP Test for Normal Mean}
Let $X_1, X_2, \ldots, X_{25}$ denote a random sample of size 25 from a normal distribution $N(\theta, 100)$. Find a uniformly most powerful critical region of size $\alpha = 0.10$ for testing $H_0: \theta = 75$ against $H_1: \theta > 75$.
\end{exercise}

\begin{exercise}{8.2.11 — MLR and UMP Test for Beta Distribution}
Let $X_1, X_2, \ldots, X_n$ be a random sample from a distribution with pdf $f(x; \theta) = \theta x^{\theta-1}$, $0 < x < 1$, zero elsewhere, where $\theta > 0$. Show the likelihood has MLR in the statistic $\prod_{i=1}^n X_i$. Use this to determine the UMP test for $H_0: \theta = \theta'$ against $H_1: \theta < \theta'$, for fixed $\theta' > 0$.
\end{exercise}

\begin{exercise}{8.3.8 — LRT for Normal Mean (Two-Sided)}
Let $X_1, X_2, \ldots, X_n$ be a random sample from the normal distribution $N(\theta, 1)$. Show that the likelihood ratio principle for testing $H_0: \theta = \theta'$, where $\theta'$ is specified, against $H_1: \theta \neq \theta'$ leads to the inequality $|\overline{x} - \theta'| \geq c$.
\end{exercise}

\begin{exercise}{8.3.11 — LRT for Two Normal Variances}
Let $X_1, \ldots, X_n$ and $Y_1, \ldots, Y_m$ be independent random samples from the distributions $N(\theta_1, \theta_3)$ and $N(\theta_2, \theta_4)$, respectively. Show that the LRT for testing $H_0: \theta_3 = \theta_4$ against $H_1: \theta_3 \neq \theta_4$ leads to the $F$-test statistic
\[
F = \frac{\sum_{i=1}^n (X_i - \bar{X})^2 / (n-1)}{\sum_{i=1}^m (Y_i - \bar{Y})^2 / (m-1)}.
\]
Under $H_0$, $F \sim F(n-1, m-1)$.
\end{exercise}

\begin{exercise}{8.4.2 — SPRT for Poisson Distribution}
Let $X$ have a Poisson distribution with mean $\theta$. Find the sequential probability ratio test for testing $H_0: \theta = 0.02$ against $H_1: \theta = 0.07$. Show that this test can be based upon the statistic $\sum_{i=1}^n X_i$. If $\alpha_a = 0.20$ and $\beta_a = 0.10$, find $c_0(n)$ and $c_1(n)$.
\end{exercise}

% === 用户问答记录 ===

